\section{Latihan}
\label{sec:latihan-bab16}

\begin{enumerate}
    \item \textbf{Konseptual.} Jelaskan mengapa \eng{disassembly}
          tidak dapat mengembalikan kode sumber asli, walaupun
          \eng{assembly} dan kode mesin bersifat satu-ke-satu.
          Sebutkan tiga informasi yang hilang selama perakitan dan
          jelaskan mengapa informasi itu tidak tersimpan di dalam
          berkas \file{.exe}.

    \item \textbf{Konseptual.} Byte \hex{48} dapat berarti instruksi
          \asm{DEC AX} atau karakter \code{'H'}. Jelaskan bagaimana
          keadaan ini menyulitkan \eng{disassembly}, lalu sebutkan
          dua cara praktis untuk menentukan batas antara daerah kode
          dan daerah data pada sebuah program.

    \item \textbf{Analisis hasil disassembly.} Perhatikan potongan
          berikut.

\begin{lstlisting}[language=MASM, numbers=none]
1D2F:0200 B90A00    MOV  CX,000A
1D2F:0203 31DB      XOR  BX,BX
1D2F:0205 0304      ADD  AX,[SI]
1D2F:0207 83C602    ADD  SI,+02
1D2F:020A E2F9      LOOP 0205
1D2F:020C 8906F401  MOV  [01F4],AX
\end{lstlisting}

          \begin{enumerate}
              \item Pola apa yang tampak pada potongan tersebut?
              \item Berapa kali instruksi \lstasm|ADD AX, [SI]|
                    dieksekusi? Jelaskan cara Anda memperolehnya.
              \item Berapa jarak relatif yang disimpan pada byte
                    \hex{E2 F9}, dan mengapa \eng{disassembler} dapat
                    menampilkan alamat \hex{0205}?
          \end{enumerate}

    \item \textbf{Analisis metrik.} Sebuah program memuat perulangan
          yang dijalankan 20 kali dengan isi
          \lstasm|ADD AX, [SI]|, \lstasm|ADD SI, 2|, dan
          \lstasm|LOOP|. Hitung perkiraan jumlah siklus seluruh
          perulangan dengan memakai Tabel~\ref{tab:siklus-8086-perkiraan}.
          Tuliskan rincian perhitungan Anda, dan sebutkan dua alasan
          mengapa angka hasil hitungan belum tentu sama dengan waktu
          sesungguhnya.

    \item \textbf{\eng{Strength reduction}.} Tuliskan hasil
          penggantian untuk potongan berikut, lalu sebutkan syarat
          yang harus dipenuhi agar penggantian itu sahih.

\begin{lstlisting}[language=MASM, numbers=none]
    MOV AX, 0
    ADD AX, nilai
    MOV BX, 16
    MUL BX
    CMP AX, 0
    JE  kosong
\end{lstlisting}

    \item \textbf{Optimasi dan trade-off.} Diberikan dua varian
          pengalian dengan empat: (a) \lstasm|MOV CL, 2| diikuti
          \lstasm|SHL AX, CL|, dan (b) dua kali \lstasm|SHL AX, 1|.
          Hitung ukuran byte dan perkiraan siklus keduanya. Varian
          mana yang Anda pilih bila program harus muat pada memori
          64 byte? Varian mana yang Anda pilih bila program berjalan
          di dalam perulangan seribu kali? Jelaskan alasan tiap
          jawaban.

    \item \textbf{Pengukuran.} Rancang pengukuran untuk
          membandingkan dua versi program yang mengerjakan pekerjaan
          sama. Sebutkan minimal empat syarat agar perbandingan itu
          adil, lalu tuliskan potongan kode yang mencatat waktu mulai
          dan waktu berhenti dengan \code{INT 1Ah} fungsi \hex{00}.
          Jelaskan mengapa percobaan pertama biasanya dibuang.

    \item \textbf{Penulisan program.} Tulislah program lengkap yang
          mengalikan setiap elemen larik \code{DW 3, 5, 7, 9, 11}
          dengan empat, lalu menjumlahkan seluruh hasil perkaliannya
          ke dalam \reg{DX}. Gunakan \asm{SHL}, bukan \asm{MUL}.
          Sertakan penghitung iterasi pada variabel tersendiri agar
          jumlah iterasi dapat diverifikasi, lalu periksa kebenaran
          hasilnya secara manual.
\end{enumerate}

\begin{refleksibox}
Jawab pertanyaan berikut pada jurnal praktikum Anda.
\begin{enumerate}
    \item Sebelum praktikum ini, apakah Anda pernah merasa yakin
          bahwa sebuah perubahan pasti mempercepat program? Setelah
          mengukur, apakah dugaan Anda terbukti? Bagian mana yang
          paling berbeda dari dugaan?
    \item Teknik optimasi mana yang paling sulit Anda terapkan
          dengan benar, dan kesalahan seperti apa yang pernah Anda
          buat saat menerapkannya?
    \item Menurut Anda, pada situasi seperti apa sebuah program
          sebaiknya dibiarkan dalam bentuk naif saja? Kaitkan
          jawaban Anda dengan nilai keterbacaan dan kemudahan
          pemeliharaan.
\end{enumerate}
\end{refleksibox}
